\section{Simplicity does not single one out}\label{sec:simplicity}

The most common preference in the search for a theory of everything is simplicity:
among the candidates that fit the data, prefer the shortest. Which candidates fit is
fixed by the record. Which of them is shortest depends on the reference machine, the
universal computer in which lengths are measured, and there are always rivals that fit
equally well, are nearly as short, and disagree about what has not been observed.

\subsection{Description length as a constraint}

\begin{definition}[Prefix description system]\label{con:def:prefix}
A \emph{prefix description system} is a partial computable function
$U:\{0,1\}^*\times\Str\rightharpoonup\Str$ whose set of programs
$D_U=\{p : U(p,x)\conv \text{ for some } x\}$ is prefix-free. A program
$p\in D_U$ determines the model $x\mapsto U(p,x)$; by the $s$-$m$-$n$ theorem an index
$e_p$ with $\ph{e_p}=U(p,\cdot)$ is computable from $p$. $U$ is \emph{universal} if
for every prefix description system $V$ there is a string $\sigma_V$ with
$U(\sigma_V q,x)\simeq V(q,x)$ for all $q,x$. Universal systems exist
\citep{chaitin1975,li-vitanyi2019}.
\end{definition}

\begin{definition}[Description length of a record]\label{con:def:KR}
For a universal $U$ and a nonempty finite functional record $\rec$, a program
$p\in D_U$ is \emph{consistent with $\rec$} if $U(p,x)=y$ for all $(x,y)\in\rec$.
Let
\[
  \KU{U}(\rec)=\min\{|p| : p\in D_U \text{ consistent with } \rec\},
\]
and let the \emph{MDL selection relative to $U$} be the set of models $e_p$ for which
$p$ is consistent with $\rec$ and $|p|=\KU{U}(\rec)$. For a string $s$ we write
$\KU{U}(s)=\KU{U}(\{(\varepsilon,s)\})$; $\K$ with no subscript denotes $\KU{U}$ for a
fixed reference $U$.
\end{definition}

$\KU{U}(\rec)$ is finite, because a lookup table for $\rec$ is a consistent program.
Minimum description length (MDL) is a constraint in the sense of
\cref{def:constraint}: let $c_U(e)=\min\{|p| : p\in D_U,\ e_p=e\}$, with
$c_U(e)=\infty$ if there is no such $p$. Since $p$ is consistent with $\rec$ exactly
when $e_p\in\Cons(\rec)$, $\Sel_{c_U}(\rec)$ is the MDL selection relative to $U$. It is
\emph{imposed} in the sense of \cref{def:constraint}: by \cref{thm:con-machine}(i)
below, the same record yields different MDL selections under different universal
machines.

\subsection{The choice of machine}

\begin{theorem}[T3: MDL is machine-relative]\label{thm:con-machine}
Let $U$ be universal, $\rec$ a nonempty finite functional record, and
$e\in\Cons(\rec)$. There is a universal prefix description system $U'$ such that
\begin{enumerate}[label=(\roman*)]
  \item for every record $\rec'\sqsupseteq\rec$ with $e\in\Cons(\rec')$, the
        MDL selection relative to $U'$ has a unique shortest consistent program,
        and it computes $\ph{e}$;
  \item $\KU{U'}(\rec')\le\KU{U}(\rec')+1$ for every nonempty functional record
        $\rec'$;
  \item $\KU{U}(\rec')\le\KU{U'}(\rec')+\KU{U}(e)+O(1)$ for every nonempty functional
        record $\rec'$, where $\KU{U}(e)$ is the length of a shortest $U$-program that
        outputs $e$ and the $O(1)$ depends only on $U$.
\end{enumerate}
\end{theorem}

\begin{proof}
Define $U'(0,x)\simeq\ph{e}(x)$ and $U'(1p,x)\simeq U(p,x)$, with $U'$ undefined on
all other programs. Then $D_{U'}\subseteq\{0\}\cup 1D_U$, which is prefix-free: $0$ is
incomparable with every string beginning with $1$, and $1p\sqsubseteq 1q$ forces
$p\sqsubseteq q$, hence $p=q$ for $p,q\in D_U$. $U'$ is universal with
$\sigma'_V=1\sigma_V$.

(i)~Since $e\in\Cons(\rec')$ and $\rec'\neq\emptyset$, $\ph{e}$ halts on some input,
so $0\in D_{U'}$ and $0$ is consistent with $\rec'$. The empty string is not in
$D_{U'}$, and $1\in D_{U'}$ only if $\varepsilon\in D_U$, which prefix-freeness and
universality exclude. So $0$ is the unique shortest consistent program.

(ii)~Each $U$-program $p$ becomes the $U'$-program $1p$.

(iii)~Let $V$ read a $U$-program $a$ (self-delimiting, since $a\in D_U$), decode
$e=U(a,\varepsilon)$, then read a $U'_e$-program $q$ and simulate $U'_e(q,\cdot)$, where
$U'_e$ is the machine $U'$ built from $e$; the simulation is uniform in $e$. Each
$U'$-program $q$ then has the $U$-program $\sigma_V a q$, where $a$ is a shortest
$U$-program for $e$, so $|a|=\KU{U}(e)$.
\end{proof}

The invariance theorem \citetext{\citealp[\S3]{kolmogorov1965}; \citealp{li-vitanyi2019}} says that two universal
machines assign description lengths within an additive constant of each other.
\Cref{thm:con-machine} says what that constant means on finite data. One extra bit in
the reference machine (ii) makes any chosen consistent candidate the unique simplest
one, and it stays so until the record refutes it (i). The price, measured in the
original machine, is at most about $\KU{U}(e)$, the length of the candidate being favored
(iii). Promoting a candidate far from the minimum, such as a lookup table for the
record, needs a change of machine with large translation constants (the converse
after \cref{thm:con-gap}). Promoting a near-tie needs only a
small one, in the following sense.

\begin{definition}[Pointwise cost of a change of machine]\label{con:def:pointwise}
Universal prefix description systems $U,U'$ differ \emph{pointwise by at most $c$} if
$|\KU{U}(\rec')-\KU{U'}(\rec')|\le c$ for every nonempty finite functional record
$\rec'$.
\end{definition}

The pointwise cost bounds how far a change of machine moves any description length.
It is not the translation constant of the invariance theorem, the length of a prefix
$\sigma$ with $U(\sigma q,x)\simeq U'(q,x)$ for all $q,x$; that constant bounds the
pointwise cost from one side, but it need not be small.

\begin{theorem}[T3, gap form: near-ties are machine-relative]\label{thm:con-gap}
Let $U$ be universal, $\rec$ a nonempty finite functional record, $\delta\ge0$ an
integer, and $q\in D_U$ a program consistent with $\rec$ with
$|q|\le\KU{U}(\rec)+\delta$. There is a universal prefix description system $U'$ such
that
\begin{enumerate}[label=(\roman*)]
  \item for every record $\rec'\sqsupseteq\rec$ with which $q$ is consistent, the
        unique shortest $U'$-program consistent with $\rec'$ computes the model
        $x\mapsto U(q,x)$;
  \item $\KU{U'}(\rec')\le\KU{U}(\rec')+1$ and $\KU{U}(\rec')\le\KU{U'}(\rec')+\delta$
        for every nonempty finite functional record $\rec'$, so $U$ and $U'$ differ
        pointwise by at most $\max(\delta,1)$.
\end{enumerate}
\end{theorem}

\begin{proof}
Let $\ell=\max(0,|q|-\delta-1)$ and $t=0^\ell$. Define $U'(0s,x)\simeq U(s,x)$ for
every $s$ and $U'(1t,x)\simeq U(q,x)$, with $U'$ undefined on all other programs. Then
$D_{U'}\subseteq 0D_U\cup\{1t\}$, which is prefix-free: $1t$ is incomparable with
every string beginning with $0$, and $0s\sqsubseteq 0s'$ forces $s=s'$ for
$s,s'\in D_U$. $U'$ is universal with $\sigma'_V=0\sigma_V$. Every $U'$-program other
than $1t$ has the form $0s$ with $s\in D_U$ and computes the same model as $s$, so
\begin{equation}\label{con:eq:gap}
  \KU{U'}(\rec')=
  \begin{cases}
    \min\bigl(\KU{U}(\rec')+1,\ \ell+1\bigr) & \text{if $q$ is consistent with $\rec'$},\\
    \KU{U}(\rec')+1 & \text{otherwise}.
  \end{cases}
\end{equation}
Note that $\KU{U}(\rec')\ge1$ for every nonempty $\rec'$: the empty program is not in
$D_U$, since otherwise $D_U=\{\varepsilon\}$ and $U$ would not be universal.

(i)~Let $\rec'\sqsupseteq\rec$ with $q$ consistent with $\rec'$. A program consistent
with $\rec'$ is consistent with $\rec$, so $\KU{U}(\rec')\ge\KU{U}(\rec)\ge|q|-\delta$.
Hence $\ell+1\le\max(1,|q|-\delta)\le\KU{U}(\rec')<\KU{U}(\rec')+1$. So $1t$, which is
consistent with $\rec'$, is strictly shorter than every consistent program $0s$, and it
computes $x\mapsto U(q,x)$.

(ii)~The first inequality is immediate from \eqref{con:eq:gap}. For the second, if
$\KU{U'}(\rec')=\KU{U}(\rec')+1$ there is nothing to prove. Otherwise
$\KU{U'}(\rec')=\ell+1$ and $q$ is consistent with $\rec'$, so
$\KU{U}(\rec')\le|q|\le\ell+1+\delta$.
\end{proof}

Two features of the theorem matter. First, the cost is pointwise only. The proof gives
no simulation of $U'$ in $U$ shorter than about $\KU{U}(q)$ bits, since $U'$ must be
told $q$; so the translation constant may be as large as the bound of
\cref{thm:con-machine}(iii), while no description length moves by more than
$\max(\delta,1)$. Second, the converse holds for translation constants. If
$U(\sigma q',x)\simeq U'(q',x)$ for all $q',x$ and $\KU{U'}\le\KU{U}+c'$, then the
shortest $U'$-program consistent with $\rec$, prefixed by $\sigma$, is a $U$-program
for the same model, so the $U'$-selection lies within $|\sigma|+c'$ bits of the
$U$-minimum. A change of machine with small translation constants reorders only
near-ties.

\begin{remark}[``Unnatural'' machines]\label{con:rem:natural}
The standard reply is that $U'$ is unnatural. But naturalness of the
reference machine is itself an imposed constraint, and no known principle privileges
one universal machine over another. The most direct attempt, a stationary distribution
for the process of universal machines simulating one another, does not exist on the
set of all universal machines; it exists only on restricted subclasses, which
reintroduces the choice \citep{muller2010}.
\end{remark}

\subsection{Near-minimal rivals}

\begin{proposition}[Near-minimal rivals]\label{thm:con-nearmin}
Let $U$ be universal, $\rec$ a nonempty finite functional record, and $p$ a shortest
program consistent with $\rec$. For every $x\notin\dom(\rec)$ and every $y\in\Str$
there is a program $q$ consistent with $\rec$ such that $U(q,x)=y$ and
\[
  |q|\le\KU{U}(\rec)+\K(x,y)+O(1),
\]
where the $O(1)$ depends only on $U$. In particular, for every such $x$ there is a
consistent $q$ that disagrees at $x$ with the shortest program $p$ (that is,
$U(q,x)\ne U(p,x)$, including the case $U(p,x)\divg$) and
$|q|\le\KU{U}(\rec)+\K(x)+O(1)$.
\end{proposition}

\begin{proof}
Let $V$ be the prefix description system that, on program $s$ and input $z$, finds the
unique prefix $a$ of $s$ in $D_U$ (by dovetailing $U$ on the prefixes of $s$), requires
$U(a,\varepsilon)=\langle x,y\rangle$, confirms that the remainder $p$ of $s$ lies in
$D_U$, and then outputs $y$ if $z=x$ and $U(p,z)$ otherwise. $D_V\subseteq\{ap : a,p\in
D_U\}$, which is prefix-free. Take $a$ a shortest $U$-program for $\langle x,y\rangle$
and set $q=\sigma_V ap$. For $(x',y')\in\rec$ we have $x'\ne x$, so
$U(q,x')=U(p,x')=y'$; and $U(q,x)=y$. For the last claim, since $|\Sig|\ge2$, choose a
single symbol $y$ different from $U(p,x)$ (any symbol if $U(p,x)\divg$); then
$\K(x,y)\le\K(x)+O(1)$.
\end{proof}

\begin{corollary}[Near-minimal rivals can be promoted]\label{con:cor:promote}
Let $U$, $\rec$ and $p$ be as in \cref{thm:con-nearmin}, and $x\notin\dom(\rec)$. There
is a universal $U'$ that differs from $U$ pointwise by at most $\K(x)+O(1)$ and whose
unique shortest program consistent with $\rec$ computes a model that disagrees with $p$
at $x$.
\end{corollary}

\begin{proof}
Take the rival $q$ of \cref{thm:con-nearmin} that disagrees with $p$ at $x$, and apply
\cref{thm:con-gap} with $\delta=|q|-\KU{U}(\rec)\le\K(x)+O(1)$.
\end{proof}

Two margins are in play here, and they are different. The first is the margin of a
\emph{change of machine}: by \cref{con:cor:promote}, a change that moves no description
length by more than $\K(x)+O(1)$ bits reverses which candidate MDL selects, in favor of
one that predicts differently at $x$. For a setting $x$ of bounded complexity this
margin is a constant. For the canonical record of length $n$ every unobserved setting
is some $x_m$ with $m\ge n$; the next one costs $\K(x_n)\le\K(n)+O(1)$, at most about $\log n$
bits, and some later ones, such as $x_{2^k}$ with $2^k\ge n$, cost only about
$\log\log n$.

The second margin is that of \emph{algorithmic statistics} for a fixed machine.
Algorithmic statistics draws the most that can be drawn from a record without a choice
among near-ties. The Kolmogorov structure function and the minimal sufficient statistic
split the record's information into a meaningful part, a simplest model for which the
record is typical, and an accidental remainder
\citep{gacs2001,vereshchagin-vitanyi2004}. Three features of that split leave the
near-ties standing. It is determined only up to additive terms logarithmic in the
length $n$ of the record; its models are sets or distributions that contain the
record, so it describes what has been observed, not how the record extends; and it is
neither computable nor usefully approximable \citep[Sec.~VII]{vereshchagin-vitanyi2004}.
The rivals of \cref{thm:con-nearmin} at the next setting cost at most
$\K(n)+O(1)$, about $\log n$, extra bits, so they lie inside the margin within which
the structure function, for any fixed machine, declares itself indifferent. This margin
does not depend on a change of machine.

\subsection{Other preferences, and serious rivals}

Simplicity is not the only preference at work in the search for a final theory.
Naturalness, beauty, symmetry, and unification are also orderings on candidate models,
and so constraints in the sense of \cref{def:constraint}: each ranks models that the
record cannot tell apart.

Nor are the surviving rivals only gerrymanders. Physics supplies rivals of comparable
compactness that agree on every prediction. When two formulations are related by a
computable, prediction-preserving translation $\tau$, their description lengths differ
by at most $\K(\tau)+O(1)$: a program for one formulation, prefixed by a
self-delimiting program for $\tau$, is a program for the other. For the explicit
translations below $\K(\tau)$ is the length of a short procedure, so if one formulation
is a shortest consistent program, the other is a near-tie of at most $\K(\tau)+O(1)$
bits, and by \cref{thm:con-gap} a change of machine that moves no description length by
more than that makes it the unique selection. Newtonian, Lagrangian and Hamiltonian
mechanics agree, for conservative systems with a regular Lagrangian, through the
Legendre transform, while their primitives differ: forces along trajectories, a
variational principle over histories, flows on phase space. Heisenberg's matrix
mechanics and Schr\"odinger's wave mechanics began from incompatible pictures and were
shown equivalent, incompletely by Schr\"odinger \citep{schrodinger1926equiv} and
rigorously by von Neumann \citep{vonneumann1932}; string-theoretic dualities are
candidates of the same kind (\cref{sec:ideal}).
